\section{更改字体}

\subsection{更改单张幻灯片上的字体}

若要更改单个段落或短语的字体，请选择要更改字体的文本。

若要更改某一占位符中所有文本的字体，请选择该占位符中的所有文本，或单击该占位符。

在“开始”选项卡的“字体”组中，从“字体”列表中选择一个字体。

\subsection{更改整个演示文稿的字体}

更改整个演示文稿的字体有以下几种方法：

\begin{enumerate}
    \item 在“设计”选项卡的“变体”组中，从“字体”列表中选择一种字体方案或者新建自定义方案。
    \item 在“视图”选项卡上的“母版视图”组中，单击“幻灯片母版”。在“幻灯片母版”选项卡的“背景”组中单击“字体”，然后从列表中选择一个字体。
    \item 在“开始”选项卡上的“编辑”组中，选择“替换”，然后选择“替换字体”。选择要替换的字体和要使用的字体，然后选择“替换”，再选择“关闭”。
\end{enumerate}

\subsection{在 PowerPoint 中更改默认字体}

可以通过创建一个自定义模板来保存上述字体设置。此模板保存字体更新，可用于将来的演示文稿。
\begin{enumerate}
    \item 根据自己的需要设置默认标题及正文字体并保存。
    \item 点击幻灯片菜单栏【设计】左侧的【主题】下拉菜单，选择【保存当前主题】；再将右键点击该自定义主题，选择【设置为默认主题】。
    \item 每次新建幻灯片，默认都是刚刚自定义的主题模版了。
\end{enumerate}

\subsection{更改幻灯片上页脚中的字体}

\begin{enumerate}
    \item 在“视图”选项卡上，单击“幻灯片母版”。
    \item 在缩略图窗格的顶部，单击幻灯片母版选中它。
    \item 在幻灯片母版上突出显示任何页脚元素（例如日期、页脚文本或幻灯片编号），然后在“开始”选项卡上，选择“字体”和“段落”组中所需的字体格式。
\end{enumerate}


\section{幻灯片母版}

\subsection{什么是幻灯片母版？}

若要使所有的幻灯片包含相同的字体和图像（如徽标），在一个位置中便可以进行这些更改，即幻灯片母版，而这些更改将应用到所有幻灯片中。 若要打开“幻灯片母版”视图，请在“视图”选项卡上选择“幻灯片母版”

母版幻灯片是窗口左侧缩略图窗格中最上方的幻灯片。 与母版版式相关的幻灯片显示在此母版幻灯片下方。

编辑幻灯片母版时，基于该母版的所有幻灯片将包含这些更改。 但是，所做的大部分更改最有可能成为与此母版相关的幻灯片版。

使用幻灯片母版统一演示文稿的的颜色、字体或效果。

\subsection{主题}

主题是相互辉映的一组颜色、字体和特殊效果（如阴影、反射、三维效果等）。 PowerPoint 提供多种设计主题 - 包括协调的配色方案、背景、字体样式和占位符的放置。我们在“普通”视图中的“设计”选项卡上向你提供这些预先设计的主题。 你还可以从 templates.office.com 获取更多主题。

您在演示文稿中使用的每个主题包括一个幻灯片母版和一组相关版式。 如果希望演示文稿包含多个主题，请再添加一个幻灯片母版，并向其应用一个主题。

\subsection{幻灯片版式}

在“幻灯片母版”视图中更改和管理幻灯片版式。 每个主题都有多个幻灯片版式。 选择与幻灯片内容最为匹配的版式；某些版式更适用于文本内容，某些版式更适用于图形内容。

每个幻灯片版式的设置各不相同 — 每个版式上的不同位置中具有不同类型的占位符。

每个幻灯片母版具有名为“标题幻灯片版式”的一个相关幻灯片版式，并且每个主题针对该版式以不同方式排列文本和其他的对象占位符，具有不同的颜色、字体和效果。

你可以更改版式的任何内容以满足你的需要。 当你更改版式，然后转到普通视图时，你之后添加的每张幻灯片都将基于此版式并反映版式更改后的外观。 但是，如果演示文稿中有基于版式旧版本的现有幻灯片，你需要对这些幻灯片重新应用版式。


\section{动画效果}

\subsection{为文本和对象添加动画}

\begin{enumerate}
    \item 选择要添加动画的对象或文本。
    \item 选择“动画”并选择一种动画。
    \item 选择“效果选项”并选择一种效果。
\end{enumerate}

要对同一对象应用多个动画效果，请选择该对象，单击“添加动画”，然后选择另一个动画效果。

\subsection{选择启动动画的方式}

\begin{enumerate}
    \item 单击时：单击幻灯片时启动动画。
    \item 与上一动画同时：与序列中的上一动画同时播放动画。
    \item 上一动画之后：上一动画出现后立即启动动画。
    \item 持续时间：延长或缩短效果。
    \item 延迟：效果运行之前增加时间。
\end{enumerate}

\subsection{一次显示一段文本}

通过名为“按段落”的动画效果选项，可使列表项一次显示一段。

\begin{enumerate}
    \item 选择包含文本的框。
    \item 选择“动画”选项卡，然后选择动画。
    \item 选择“效果选项”，然后选择“按段落”，使文本段落一次显示一个段落。
\end{enumerate}

默认情况下，在“幻灯片放映”中演示时，每个段落都将显示以响应单击。 这样，您可以控制每个段落的显示时间。 可以使用功能区“动画”选项卡最右端上的“开始”、“持续时间”和“延迟”控件修改此设置。

\subsection{使文本一次显示一个字母}

您还可以通过使段落中的字符一次显示一个来创建“键入”视觉效果。

\begin{enumerate}
    \item 在幻灯片上，选择包含文本的框。
    \item 选择“动画”选项卡，然后选择动画。
    \item 选择“动画窗格”，将打开“动画窗格”窗口。
    \item 在“动画窗格”中，选择动画旁边的箭头，然后选择“效果选项”。
    \item 在对话框中的“效果”选项卡上的“增强功能”下，选择“将文本制作为动画”旁边的箭头，然后选择“按字母”。 然后，可以在“字母之间的延迟秒数”框中更改延迟时间。
\end{enumerate}


\subsection{应用动作路径动画}

\begin{enumerate}
    \item 单击要应用动作的对象。
    \item 在“动画”选项卡中，单击“添加动画”。
    \item 向下滚动到“动作路径”下，选择一个。
    \item 如果选择 “自定义路径”选项，则会绘制希望对象采用的路径。若要停止绘制自定义路径，请按 Esc。
    \item 要执行某些操作（例如，更改动作路径的方向，编辑动作路径的各个点，锁定（以使其他人无法更改动画）或取消锁定动画），请单击“效果选项”。
    \item 
\end{enumerate}


\subsection{在PowerPoint中使用平滑切换}

平滑切换功能可将从一张幻灯片到另一张幻灯片的平滑移动具有动画效果。 